%% =================================================================
%% Wei, Mei, Zhao, Xiao, Sun, Zhang, Guo.
%% "Nondestructively identifying the mechanical behavior of soft
%%  tissues using surface deformation with an explicit inverse
%%  approach."
%% Appl. Math. Modelling 134 (2024) 126--147.
%%
%% Reconstruction of page 2 (printed p. 127) as study notes.
%% Prose: extracted verbatim from the PDF text layer via pdftotext
%%        (hyphens and line breaks cleaned up only).
%% Math:  hand-transcribed from the rendered page image to preserve
%%        equation layout that pdftotext mangles.
%% =================================================================

\documentclass[11pt]{article}

\usepackage[margin=1in]{geometry}
\usepackage{amsmath, amssymb}
\usepackage{mathrsfs}
\usepackage{fancyhdr}

\pagestyle{fancy}
\fancyhf{}
\lhead{\small Y.~Mei et al.}
\rhead{\small Applied Mathematical Modelling 134 (2024) 126--147}
\cfoot{\small 127 $\to$ \arabic{page}}
\renewcommand{\headrulewidth}{0.4pt}

\setcounter{page}{1}

\begin{document}

\noindent
\ldots{}[9--11], dietary habits [12], and the presence of pathological
conditions [13,14]. Therefore, understanding the nonhomogeneous
biomechanical behavior of these tissues is of paramount importance,
particularly in a non-destructive manner. Such knowledge can inform
medical diagnostics and treatments, potentially leading to better
patient outcomes.

Considerable effort has been devoted into research aimed at
identifying the nonhomogeneous biomechanical properties of biological
tissues. Studies [15--18] have demonstrated the feasibility of
estimating the mechanical properties of each layer in multi-layered
biological tissues, provided that the thickness of each layer is known
beforehand. However, in practice, determining the thickness of each
layer prior to addressing the identification problem can be a
challenging task. Furthermore, the interface between neighboring
layers is seldom flat, a characteristic particularly pronounced in
biological tissues. As such, a comprehensive solution requires
knowledge of both the topology of the interface between the layers
and the biomechanical properties of each layer. Additionally, the
detection of tumors embedded within normal tissues introduces another
layer of complexity. To accurately diagnose and plan treatment for
these conditions, it is crucial to determine not only the elastic
modulus values of these tumors but also their topology [19].

A general strategy for tackling this challenge involves identifying
the spatial variation of material properties, a process that
necessitates solving a large-scale optimization problem to optimize
the elastic properties of every pixel in the domain of interest,
which can be referred to as the implicit inverse method [20--23]. The
implicit inverse method has been successfully employed in tumor
detection using real clinical data [8,24]. However, this implicit
inverse method has its limitations. It not only demands significant
computational resources but is also prone to severe uniqueness
issues, which can hamper the reliability of the results. Given that
biological tissues are often known a priori to exhibit either layered
structures or inclusion-embedded structures, it's feasible to
explicitly and simultaneously identify the topology of each
nonhomogeneous region and its associated material properties. This
approach bypasses some of the computational hurdles of the implicit
method. Another inherent drawback of the implicit inverse method lies
in the quality of the reconstruction results, which can be
compromised if the full-field displacement field cannot be measured
with high accuracy. Studies [25,26] that employed the implicit
inverse method to reconstruct the nonhomogeneous elastic property
distribution of solids using surface deformation have demonstrated
this issue. While these methods are capable of visualizing the
nonhomogeneous region, the quality of the reconstruction is highly
sensitive to noise. This sensitivity arises from a disparity between
the large number of unknown optimization variables in the implicit
inverse method and the sparsity of the deformation information on the
surface. In contrast, the explicit inverse method significantly
reduces the total number of optimization variables [27--29],
potentially improving the quality of the reconstructed results. This
strategy holds promise for enhancing the accuracy of nonhomogeneous
biomechanical property identification in biological tissues.

This paper will present an explicit inverse method that utilizes
surface deformation measurements to identify the mechanical behavior
of two common nonhomogeneous structures: layered structures and
inclusion embedded structures. The organization of this work is as
follows: In the \textbf{Methods} Section, we provide a detailed
explanation of the mathematical background underlying the explicit
inverse approach. The \textbf{Results} Section presents a comparative
study between the explicit and implicit inverse approaches, using
various numerical examples. We discuss the obtained results in the
\textbf{Discussion} Section and conclude the paper with a summary in
the \textbf{Conclusion} Section.

\section*{Methods}

\subsection*{\textit{Forward problem}}

\noindent
Given the low elastic modulus of soft tissues, this study will
consider finite deformation in the analysis. The strong form of the
nonlinear elastic problem is stated as follows:
\begin{equation}
\left\{
\begin{aligned}
\mathrm{Div}(\mathbf{P}) &= \mathbf{0} && \text{in } \Omega \\
\mathbf{u}               &= \mathbf{g} && \text{on } \Gamma_g \\
\mathbf{P}\cdot\mathbf{N} &= \mathbf{T} && \text{on } \Gamma_T
\end{aligned}
\right.
\end{equation}
where $\mathbf{P} = \mathbf{F}\mathbf{S}$ is the first
Piola--Kirchhoff stress tensor where $\mathbf{F}$ is the deformation
gradient and $\mathbf{S}$ is the second Piola--Kirchhoff stress
tensor. $\Omega$ is the problem domain at the reference
configuration. $\mathbf{g}$ is the prescribed displacement on the
displacement boundary $\Gamma_g$. $\mathbf{T}$ is the prescribed
traction on the traction boundary $\Gamma_T$. The following
restrictions apply to $\Gamma_g$ and $\Gamma_T$: $\Gamma_g \cup
\Gamma_T = \partial\Omega$ and $\Gamma_g \cap \Gamma_T = \phi$. We
assume the soft tissue of interest obeys the incompressible and
neo-Hookean law:
\begin{equation}
\mathbf{S} = \mu\!\left(\mathbf{I} - \mathbf{C}^{-1}\right)
             + p\,\mathbf{C}^{-1}
\end{equation}
where $\mathbf{C} = \mathbf{F}^{\mathsf{T}}\mathbf{F}$ is the
Cauchy--Green tensor. $p$ is the hydrostatic pressure.

The weak form of the nonlinear elastic forward problem is stated as:
Find $[\mathbf{u}^h, p^h] \in \mathscr{M}^h \times \mathscr{P}^h$
such that
\begin{equation}
\mathscr{A}\!\left(\mathbf{w}^h,\mathbf{u}^h,q^h,p^h\right)
\;=\;
\int_{\Gamma_T}\! \mathbf{w}^h\!\cdot\mathbf{T}\,\mathrm{d}\Gamma,
\qquad
\forall\,[\mathbf{w}^h, q^h] \in \ell^h \times \mathscr{P}^h
\end{equation}
\begin{equation}
\mathscr{A}\!\left(\mathbf{w}^h,\mathbf{u}^h,q^h,p^h\right)
\;\equiv\;
\int_\Omega \nabla\mathbf{w}:\mathbf{P}\,\mathrm{d}\Omega
\;+\;
\int_\Omega q\,(J-1)\,\mathrm{d}\Omega
\;-\;
\sum_{i=1}^{NE}\!
 \Bigl(\tau\,\mathrm{div}(\mathbf{F}^h\mathbf{S}^h),\;
       (\mathbf{F}^h)^{-\mathsf{T}}\nabla q^h\Bigr)_{\Omega_i}
\end{equation}

\end{document}
